\section*{NeurIPS Paper Checklist}

\begin{enumerate}

\item {\bf Claims}
\item[] Question: Do the main claims in the abstract and introduction accurately
reflect the paper's contributions and scope?
\item[] Answer: \answerYes{}
\item[] Justification: The abstract states three quantitative claims (6.13-point
gap, 9,000 questions tested, VisFlow gain reduces to 2.14 points). All three
are supported by the data in Tables~2, 3, and~4.
\item[] Guidelines:
\begin{itemize}
    \item The answer NA means that the abstract and introduction do not include
    the claims made in the paper.
    \item The abstract and/or introduction should clearly state the claims made,
    including the key metrics and datasets used.
    \item The claims made should match theoretical and experimental results, and
    reflect how much the results can be expected to generalize to other settings.
    \item It is fine to include aspirational goals as motivation as long as it is
    clear that these goals are not attained in the submitted work.
\end{itemize}

\item {\bf Limitations}
\item[] Question: Does the paper discuss the limitations of the work?
\item[] Answer: \answerYes{}
\item[] Justification: Section~\ref{sec:limitations} discusses scope (tested
only LLaVA-1.5-7B), the descriptive nature of the 100-question diagnostic,
and the inability of inference-time methods to improve over the corrected
adversarial baseline.
\item[] Guidelines:
\begin{itemize}
    \item The answer NA means that the paper has no limitation while the answer
    No means that the paper has limitations, but those are not discussed in
    the paper.
    \item The authors are encouraged to create a separate ``Limitations'' section
    in their paper.
    \item The paper should point out any strong assumptions and how robust the
    results are to violations of these assumptions (e.g., independence
    assumptions, scalability assumptions, IID assumptions, etc.).
    \item The limitations of the work should be clearly explained.
    \item Authors should think about how their results would be affected if the
    assumptions were violated.
    \item There might be papers or assumptions that are not stated in the paper
    but are implicit in the code/experiments.
    \item Authors should reflect on how these assumptions might be violated in
    practice and what the consequences would be.
    \item Authors should reflect on the scope of the claims, e.g., if the method
    was only tested on a few datasets or with a few runs, a claim about
    large-scale performance is unjustified.
\end{itemize}

\item {\bf Theory Assumptions and Proofs}
\item[] Question: For each theoretical result, does the paper provide the full
set of assumptions and a complete proof?
\item[] Answer: \answerNA{}
\item[] Justification: This paper makes no theoretical claims. All contributions
are empirical.
\item[] Guidelines:
\begin{itemize}
    \item The answer NA means that the paper does not include theoretical results.
    \item All the theorems, formulas, and proofs in the paper should be clearly
    numbered and cross-referenced.
    \item It would be recommended to include the proof in the main paper, but
    when the proof is lengthy, it can be included in the appendix.
    \item Proofs of central theorems should be included in the main paper.
    \item Note that it is acceptable to include proofs of some results outside
    the scope of the main claim, e.g., as supporting material.
    \item Authors should be sure to check whether they have included all the
    necessary assumptions.
\end{itemize}

\item {\bf Experimental Result Reproducibility}
\item[] Question: Does the paper fully disclose all the information needed to
reproduce the main experimental results of the paper to the extent that it
affects the main claims and/or conclusions of the paper (regardless of whether
the code and data are provided)?
\item[] Answer: \answerYes{}
\item[] Justification: Section~\ref{sec:artifact} describes the model, hardware,
quantization settings, generation parameters, and evaluation protocol in full.
The eight token IDs and the decision rule are stated explicitly. The evaluation
script is released on GitHub.
\item[] Guidelines:
\begin{itemize}
    \item The answer NA means that the paper does not include experiments.
    \item If the paper includes experiments, a No answer to this question will
    not be perceived well by the reviewers.
    \item The experimental setting should be presented in the main paper; more
    extensive details can be provided in the appendix.
    \item Depending on the contribution, reproducibility can be accomplished in
    different ways.
    \item Authors are strongly encouraged to provide source code and data.
\end{itemize}

\item {\bf Open Access to Data and Code}
\item[] Question: Does the paper provide open access to the data and code used
to produce the main experimental results in the paper?
\item[] Answer: \answerYes{}
\item[] Justification: The evaluation script, diagnostic tool, and all 9,000
prediction records (predictions, labels, token IDs, and logits for all 16
yes/no token IDs) are released at the GitHub repository listed in
Section~\ref{sec:intro}. POPE uses COCO images, which are publicly available
under the COCO Terms of Use.
\item[] Guidelines:
\begin{itemize}
    \item The answer NA means that the paper does not include experiments
    requiring code.
    \item Please see the NeurIPS code and data submission guidelines
    (\url{https://nips.cc/public/guides/CodeSubmissionPolicy}) for more details.
    \item While we encourage the release of code and data, we understand that
    this might not be possible, so ``No'' is an acceptable answer. Papers
    cannot be rejected simply for not including code, unless this is central
    to the paper's claims.
    \item The instructions should contain the exact command(s) needed to run to
    reproduce the results shown in the paper.
    \item The authors should make sure results are reproducible from the code
    they submit.
    \item In the submission phase, to preserve anonymity, the authors should
    release anonymized versions of their code and data. Once deanonymized, the
    authors are free to release the data under the license of their choosing.
\end{itemize}

\item {\bf Experimental Setting/Details}
\item[] Question: Does the paper specify all the training and test details
(e.g., architecture, data splits, hyperparameters, how they were chosen)
necessary to understand the results?
\item[] Answer: \answerYes{}
\item[] Justification: No training was performed. Evaluation hyperparameters
(quantization, generation parameters) are specified in
Section~\ref{sec:artifact}. Inference-time correction hyperparameters
(beta values, threshold values) are listed in Section~\ref{sec:correction}.
\item[] Guidelines:
\begin{itemize}
    \item The answer NA means that the paper does not include experiments.
    \item The experimental details should be included in the paper or, if
    lengthy, as an appendix.
    \item They should be reported regardless of whether the algorithm is
    included in a contribution.
    \item Depending on the contribution, these details might include:
    \begin{itemize}
        \item Model and architecture details
        \item Optimizer and training details
        \item Dataset details
        \item Evaluation metrics
        \item Number of evaluation runs
        \item Ablation study settings
    \end{itemize}
\end{itemize}

\item {\bf Experiment Statistical Significance}
\item[] Question: Does the paper report error bars suitably and correctly
defined or other appropriate information about the statistical significance
of the experiments?
\item[] Answer: \answerNo{}
\item[] Justification: All results use deterministic evaluation (fixed model,
no temperature sampling, full 3,000-question splits). There is no randomness
to characterize with error bars.
\item[] Guidelines:
\begin{itemize}
    \item The answer NA means that the paper does not include computational
    experiments.
    \item The authors should answer ``Yes'' if the results are accompanied by
    error bars, confidence intervals, or statistical significance tests, at
    least for the experiments that support the main claims of the paper.
    \item The factors of variability that the authors are expected to take into
    account include:
    \begin{itemize}
        \item The random seed and/or the random selection of a subset of samples
        used for the proof of concept.
        \item The batch size when the gradient computation is stochastic.
        \item The values of the hyperparameters, especially when there is a wide
        variety of possible values for the hyperparameters.
        \item The number of runs used to estimate the accuracy of the results.
    \end{itemize}
    \item A systematic description of the `degrees of freedom' in the
    experimental setup is needed even if the paper does not include
    error bars.
\end{itemize}

\item {\bf Experiments Compute Resources}
\item[] Question: For each experiment, does the paper provide sufficient
information on the compute resources (type of computinginfrastructure,
memory, time of experiments) used to reproduce those experiments?
\item[] Answer: \answerYes{}
\item[] Justification: Section~\ref{sec:artifact} reports the GPU (RTX~3070~Ti,
8~GB VRAM), quantization method, and total evaluation time (76.3 minutes for
9,000 questions). Section~\ref{sec:discussion} compares per-question times for
logit-based versus string-based evaluation.
\item[] Guidelines:
\begin{itemize}
    \item The answer NA means that the paper does not include experiments.
    \item The paper should indicate the type of compute workers CPU, GPU, or TPU,
    the memory available, and the time needed to reproduce the results (in
    combination with the other factors described above, such as
    computation of loss and gradient).
    \item For models trained from scratch, the paper should provide the number of
    GPUs used, the amount of time, and the type of GPUs used, matching the
    amount of memory available.
    \item For pre-trained models, the paper should provide the number of GPUs used
    and the amount of time to obtain the reported results.
    \item For training runs, the authors should save checkpoints routinely, in
    case of failed or interrupted experiments, and let the temporary access
    to existing checkpoints be available to the reviewers.
    \item Researchers comparing to this method should take into account both the
    temporary and long-term compute savings from the reuse of pretrained
    models.
\end{itemize}

\item {\bf Code Of Ethics}
\item[] Question: Does the research conducted in the paper conform, in every
way, with the NeurIPS Code of Ethics
\url{https://neurips.cc/public/EthicsGuidelines}?
\item[] Answer: \answerYes{}
\item[] Justification: The paper reports an evaluation artifact in a published
benchmark. It does not involve human subjects, private data, or dual-use
risks.
\item[] Guidelines:
\begin{itemize}
    \item The answer NA means that the authors have not reviewed the NeurIPS
    Code of Ethics.
    \item If the authors answer No, they should explain the special circumstances
    that require a deviation from the Code of Ethics.
    \item The authors should make sure to preserve anonymity (e.g., if there is
    a special consideration due to laws or regulations in their jurisdiction).
\end{itemize}

\item {\bf Broader Impacts}
\item[] Question: Does the paper discuss both potential positive societal impacts
and negative societal impacts of the work performed?
\item[] Answer: \answerNA{}
\item[] Justification: This paper corrects a measurement error in an existing
benchmark. It has no foreseeable negative societal impact.
\item[] Guidelines:
\begin{itemize}
    \item The answer NA means that there is no societal impact of the work
    performed.
    \item If the authors answer NA or No, they should explain why their work
    has no societal impact or why the paper does not address societal impact.
    \item Examples of negative societal impacts include potential malicious or
    unintended uses (e.g., disinformation, generating fake profiles, surveillance),
    fairness considerations (e.g., deployment of technologies that could make
    decisions that unfair to some groups), privacy considerations, and
    security considerations.
    \item The Conference does not expect papers to have a {\it comprehensive}
    discussion of potential negative societal impacts.
\end{itemize}

\item {\bf Safeguards}
\item[] Question: Does the paper describe safeguards that have been put in place
for responsible release of data or models that have potential for misuse (e.g.,
high risk data or models)?
\item[] Answer: \answerNA{}
\item[] Justification: We release only an evaluation script and prediction logs.
Neither poses misuse risk.

\item {\bf Licenses for Existing Assets}
\item[] Question: Are the creators or original owners of assets (e.g., code,
data, models) used in the paper, properly attributed and are the license terms
respected?
\item[] Answer: \answerYes{}
\item[] Justification: We cite LLaVA-1.5, POPE, COCO, CLIP, and BitsAndBytes.
COCO is used under its published Terms of Use. LLaVA-1.5 and CLIP model weights
are used under their respective licenses (Apache 2.0 and MIT).

\item {\bf New Assets}
\item[] Question: Are new assets introduced in the paper clearly documented and
is the documentation provided alongside the assets?
\item[] Answer: \answerYes{}
\item[] Justification: The evaluation script and diagnostic tool are documented
in the repository README with example commands (see Appendix~\ref{app:diagnostic}).

\item {\bf Crowdsourcing and Research with Human Subjects}
\item[] Question: For crowdsourcing experiments and research with human subjects,
does the paper include the working instructions, a description of the
compensation paid to workers, and the number of annotators?
\item[] Answer: \answerNA{}
\item[] Justification: No human subjects or crowdsourcing was involved.

\item {\bf Institutional Review Board (IRB) Approvals or Equivalent for
Research with Human Subjects}
\item[] Question: Does the paper describe potential risks incurred by study
participants as well as how these risks were mitigated, and whether and how
IRB approval (or an equivalent approval/review based on the requirements of
your country or institution) was obtained?
\item[] Answer: \answerNA{}
\item[] Justification: No human subjects were involved.

\end{enumerate}
